\subsection{群作用于空间的情形}

设 X 是局部紧空间，局部紧群 G 在其上连续地左作用，作用由

\[
(s, x) \mapsto s \cdot x.
\]

若 \(\mu_{1},\ldots,\mu_{n}\) 是 G 上的测度，\(\nu\) 是 X 上的测度，则称这些测度可卷积，是指它们对于映射 \((s_{1},\ldots,s_{n},x)\mapsto s_{1}\cdots s_{n}x\)，即从 \(G^{n}\times X\) 到 X 的映射，可卷积；它们的卷积积应按此映射的意义理解。

若 \(s\in \mathbf{G}\) 且 \(x\in \mathbf{X}\)，则

\[
\varepsilon_ {s} * \varepsilon_ {x} = \varepsilon_ {s x}.\tag{4}
\]

若 \(s \in \mathbf{G}\) 且 \(\mu \in \mathcal{M}(\mathbf{X})\)，则

\[
\varepsilon_ {s} * \mu = \gamma (s) \mu\tag{5}
\]

根据 §1 No.~1 的例 3。

命题 8. --- 设 \(\mu\) 是 G 上的测度，\(\nu\) 是 X 上的测度。

\begin{enumerate}
\def\labelenumi{(\roman{enumi})}
\item
  若 \(\mu\) 具有紧支集，则 \(\mu\) 与 \(\nu\) 可卷积。
\item
  若 \(\nu\) 具有紧支集，且 \(G\) 在 \(X\) 上适当作用，则 \(\mu\) 与 \(\nu\) 可卷积。
\end{enumerate}

这由 §1 No.~4 的命题 4 得出。

命题 9. --- 关于卷积，\(\mathcal{M}^{1}(\mathrm{X})\) 是 \(\mathcal{M}^{1}(\mathrm{G})\) 上的左模，而 \(\mathcal{M}(\mathrm{X})\) 和 \(\mathcal{C}'(\mathrm{X})\) 是 \(\mathcal{C}'(\mathrm{G})\) 上的左模。

这由命题 8 以及 §1 的命题 1、3 和命题 5 的推论得出。

命题 10. --- 设 \(\mu\) 是 \(G\) 上的测度，\(\nu\) 是 \(X\) 上的测度，且 \(\mu\) 与 \(\nu\) 可卷积。再假定存在正测度 \(\beta\)，定义在 \(X\) 上，使得 \(\gamma(s)\nu\) 以 \(\beta\) 为基，对每个 \(s \in G\) 都成立。则 \(\mu * \nu\) 以 \(\beta\) 为基。

令 K 是 X 的一个 \(\beta\)-零测紧子集。则 K 都是 \(\pmb{\gamma}(s)|\nu|\)-零测的，对每个 \(s \in G\) 均如此。现在，

\[
| \mu | * | \nu | = \int_ {\mathrm{G}} (\varepsilon_ {s} * | \nu |) d | \mu | (s)
\]

（§1 No.~5 命题 7），且映射 \(s \mapsto \varepsilon_s * |\nu|\) 弱连续（§2 命题 6）。因此根据第 V 章 §3 No.~3 定理 1，K 是 \(|\mu| * |\nu|\)-零测的。因此 \(|\mu| * |\nu|\) 以 \(\beta\) 为基（第 V 章 §5 No.~5 定理 2）。
% semantic-fragments: legacy-input-seam
\relax{}
